\RequirePackage[T1]{fontenc}
\documentclass[10pt,conference]{IEEEtran}
\usepackage{amsmath,amssymb,amsthm,graphicx,booktabs,array,url,flushend}
\usepackage[hidelinks]{hyperref}
\usepackage{tikz}
\usetikzlibrary{arrows.meta,positioning}
\emergencystretch=2em
\newtheorem{theorem}{Theorem}
\newtheorem{lemma}{Lemma}
\newtheorem{proposition}{Proposition}
\theoremstyle{definition}
\newtheorem{definition}{Definition}
\newtheorem{remarknum}{Remark}
\newcommand{\code}[1]{\texttt{#1}}
\newcommand{\G}[1]{G#1}
\newcommand{\tool}{ReleaseWeave}
\newcommand{\lang}{RIL}
\begin{document}
\title{Recovering Versioned Releases Across Build and Publish Boundaries}
\maketitle
\begin{abstract}
An npm package can build from local dependency output, while an isolated
registry-based Docker builder needs its base image published first. After
partial failure, a provider's success can also remove the reason for its
consumer's unfinished release. Our protocol combines stage-aware scheduling
with consumer-based admission: Git history determines plans, and each
consumer's completion record discharges its own obligation. The single-history
model gives deterministic planning and eight conditional recovery guarantees.
A polyrepo extension retains repository-owned records and source-revision
evidence. Any peer can start CI in either topology; deployments may centralize
workflow entry in one repository or distribute entry points across initiating
repositories. For $q$ repositories, both star and minimal-tree topologies use
$q-1$ links; a star bounds evidence routes at two hops, versus $q-1$ in the
worst-case tree. These are structural bounds, not measured fleet speedups.
In 120 controlled fault episodes spanning five dependency graphs, two release
channels, and three failure boundaries, every episode recovers after repair
without losing pending obligations. An analytical alternative that discharges
intent when its source completes loses unfinished work in 84 of the same 120
episodes (70\%). This compares admission rules, not competing tools or
production failure rates. Separate measurements on
an earlier binary, with 200 intent units, show median planning time and peak
memory rising from 0.25\,s and 21\,MiB at 1,000 packages to 19.27\,s and
92\,MiB at 10,000 packages. These controlled results support the recovery
mechanism; full implementation refinement, fleet validation, and external
exactly-once effects remain outside the claims.
\end{abstract}
\begin{IEEEkeywords}
software release engineering, continuous integration and delivery,
dependency scheduling, recovery, idempotence
\end{IEEEkeywords}
\input{introduction}
\input{system}
\input{model}
\input{guarantees}
\input{evaluation}
\input{discussion}
\clearpage
\IEEEtriggeratref{34}
\input{references}
\end{document}
